\subsection{One-Sided Deterministic Updates Can Have an Exponential Gap}
\label{app:indicatoraddress}
Large global gaps do not require complementary classifiers. Let $A$ have
$\kappa\ge1$ address bits, let $D=2^\kappa$, and write $Z_j\in\{0,1\}$
for its data bits. The old rule always predicts $-1$, and the new rule
predicts $+1$ exactly when $Z_A=1$. Their contrast is $d=Z_A\in\{0,1\}$.
\begin{proposition}\label{prop:indicatoraddress}
For this contrast, $m=\kappa+D$, $b=\kappa+1$, and
\[
 V_{\rm branch}(d,b)=1,\qquad
 V_{\rm fixed}(d,b)=\frac{D^2}{4(D-1)}.
\]
The ratio is asymptotic to $2^{b-3}$. Multiplying the contrast by any
$\epsilon\in(0,1]$ leaves the ratio unchanged.
\end{proposition}
\begin{proof}
Put $\pi=D/[2(D-1)]$. Choose a data index $J$ uniformly, query all address
bits and $Z_J$, and return the feature score
\[
 H=\pi+D\mathbf1\{J=A\}(Z_J-\pi).
\]
Its conditional mean is $d$. If $J=A$, its value is $D(d-1/2)$; otherwise
it is $\pi$. Hence its conditional second moment is
$D/4+(D-1)\pi^2/D=D^2/[4(D-1)]$ for either value of $d$.

For the converse, let $\mu$ have a uniform address and independent
Bernoulli$(\pi)$ data bits, and take $h=\pi+D(d-\pi)$.
A fixed form with $s$ address bits and $r$ data bits has
\[
 \|\E_\mu[h\mid X_T]\|_\mu^2
 =\pi^2+\pi(1-\pi)r2^s
 \le\pi^2+D\pi(1-\pi)=:F,
\]
because $r+s\le\kappa+1$ implies $r2^s\le D$.
Also $\langle d,h\rangle_\mu=D\pi-(D-1)\pi^2=F$.
The variance dual therefore gives $V_{\rm fixed}\ge F=D^2/[4(D-1)]$.
This calculation includes $D=2$, where $\pi=1$ and the lower bound is
simply one. Branching evaluates $d$ exactly, and Jensen gives its matching
coefficient one. Both coefficients scale by $\epsilon^2$.
\end{proof}

This is a constructed model family, not a claim about how frequently
trained classifiers have such gaps. It separates the structural question
from model-family labels: a GBDT or neural-network name alone determines
neither identifiability nor the global efficiency ratio.
